Early work on world models showed that compact latent dynamics can capture environment structure and support planning~\citep{ha2018world}.
Recent video systems scale spatiotemporal generation and simulation, including Sora~\citep{openai2024sora}, Veo~\citep{deepmind_veo_2025}, and Movie Gen~\citep{polyak2024movie}.
Together they suggest that high‑capacity video generators can form the basis for interface modeling when aligned with inputs and actions.

Neural‑OS explores neural desktop modeling with recurrent components and latent diffusion~\citep{rivard2025neuralos,hochreiter1997long,rombach2022high}, pointing toward OS‑level simulators.
Our position differs: we pursue a pixel‑grounded, action‑conditioned framework instantiated as interface-specific NCs for CLI and GUI without hand‑coded planners.

Structured, DOM‑based agents enable reliable web interaction and evaluation~\citep{deng2023mind2web} and are widely used in commercial assistants with LLM planning and tool use.
We treat DOM/accessibility as optional auxiliaries; a unified pixel interface with a latent runtime state generalizes beyond web UIs and keeps I/O on one timeline.

Pixel‑acting agents with multimodal backbones are increasingly capable.
Expressing control as conditioning—rather than as a separate controller—reduces brittleness and improves transfer across interfaces.

Joint-embedding predictive objectives and latent diffusion backbones provide practical building blocks~\citep{assran2023self,rombach2022high}. We align these components around a shared latent-dynamics recipe and add rigorous, pixel-grounded evaluation focusing on action–frame alignment, temporal context, and capacity. Our contribution is a pixel-based system spanning CLI/GUI interfaces without hand-coded controllers, with action-conditioned training and evaluation that measures on-screen performance (pointer accuracy, latency, \LPIPS/\FVD) and maps design choices to functional improvements.
